\section{Тестирование приложения}

\subsection{Общее описание}

В результате выполнения курсовой работы было разработано приложение телефонного справочника с графическим интерфейсом. Программа реализует все требования задания:

\begin{itemize}
    \item хранение контактов с полями: фамилия, имя, отчество, адрес, дата рождения, email, телефоны;
    \item валидация вводимых данных с помощью регулярных выражений;
    \item добавление, редактирование и удаление контактов;
    \item поиск по различным полям;
    \item сортировка по столбцам с трёхступенчатым переключением;
    \item сохранение в бинарный формат и экспорт/импорт в YAML.
\end{itemize}

\subsection{Главное окно приложения}

При запуске программы отображается главное окно с пустой таблицей контактов (рис.~\ref{fig:Screenshot_2025-12-07_at_20.44.35.png}). В верхней части расположена панель поиска с выбором поля для фильтрации. Внизу находятся кнопки управления: «Добавить», «Редактировать», «Удалить» слева и «Импорт YAML», «Экспорт YAML» справа.

\screenshot{Screenshot_2025-12-07_at_20.44.35.png}{Главное окно приложения при первом запуске}

\subsection{Добавление нового контакта}

При нажатии кнопки «Добавить» открывается диалог создания контакта (рис.~\ref{fig:Screenshot_2025-12-07_at_20.45.12.png}). Форма разделена на три группы: ФИО, контактные данные и телефоны. Обязательные поля отмечены звёздочкой.

\screenshot{Screenshot_2025-12-07_at_20.45.12.png}{Диалог создания нового контакта}

\subsection{Валидация вводимых данных}

При попытке сохранить контакт с некорректными данными отображаются сообщения об ошибках под соответствующими полями (рис.~\ref{fig:Screenshot_2025-12-07_at_20.45.30.png}). На скриншоте видны ошибки:
\begin{itemize}
    \item фамилия, имя и отчество не начинаются с заглавной буквы;
    \item email имеет неверный формат;
    \item не добавлен ни один телефон.
\end{itemize}

\screenshot{Screenshot_2025-12-07_at_20.45.30.png}{Отображение ошибок валидации}

\subsection{Корректное заполнение формы}

После исправления ошибок форма заполнена корректно (рис.~\ref{fig:Screenshot_2025-12-07_at_20.48.51.png}). Телефон добавлен в таблицу и отображается в отформатированном виде.

\screenshot{Screenshot_2025-12-07_at_20.48.51.png}{Корректно заполненная форма контакта}

\subsection{Отображение контакта в таблице}

После сохранения контакт появляется в главной таблице (рис.~\ref{fig:Screenshot_2025-12-07_at_20.49.01.png}). Телефон отображается в международном формате, дата~--- в формате ДД.ММ.ГГГГ.

\screenshot{Screenshot_2025-12-07_at_20.49.01.png}{Главное окно с одним контактом}

\subsection{Работа с несколькими контактами}

После добавления нескольких контактов таблица заполняется данными (рис.~\ref{fig:Screenshot_2025-12-07_at_20.51.22.png}, рис.~\ref{fig:Screenshot_2025-12-07_at_20.53.41.png}). При выборе строки активируются кнопки «Редактировать» и «Удалить».

\screenshot{Screenshot_2025-12-07_at_20.51.22.png}{Таблица с двумя контактами}

\screenshot{Screenshot_2025-12-07_at_20.53.41.png}{Таблица с тремя контактами}

\subsection{Редактирование контакта}

При двойном клике по строке или нажатии кнопки «Редактировать» открывается диалог с заполненными данными выбранного контакта (рис.~\ref{fig:Screenshot_2025-12-07_at_20.54.17.png}). У контакта Иванов добавлено отчество и второй телефон рабочего типа.

\screenshot{Screenshot_2025-12-07_at_20.54.17.png}{Диалог редактирования контакта}

\subsection{Удаление контакта}

При нажатии кнопки «Удалить» появляется диалог подтверждения с именем удаляемого контакта (рис.~\ref{fig:Screenshot_2025-12-07_at_20.54.57.png}).

\screenshot{Screenshot_2025-12-07_at_20.54.57.png}{Диалог подтверждения удаления}

После подтверждения контакт удаляется из таблицы (рис.~\ref{fig:Screenshot_2025-12-07_at_20.55.06.png}).

\screenshot{Screenshot_2025-12-07_at_20.55.06.png}{Таблица после удаления контакта}

\subsection{Добавление контакта после удаления}

Демонстрация добавления нового контакта после удаления предыдущего (рис.~\ref{fig:Screenshot_2025-12-07_at_20.56.56.png}).

\screenshot{Screenshot_2025-12-07_at_20.56.56.png}{Таблица после добавления нового контакта}

\subsection{Экспорт в YAML}

При нажатии кнопки «Экспорт YAML» открывается стандартный диалог сохранения файла (рис.~\ref{fig:Screenshot_2025-12-07_at_20.57.14.png}). По умолчанию предлагается имя \texttt{phonebook.yaml} и фильтр по расширениям \texttt{*.yaml} и \texttt{*.yml}.

\screenshot{Screenshot_2025-12-07_at_20.57.14.png}{Диалог экспорта в YAML}

Можно указать путь к папке проекта напрямую (рис.~\ref{fig:Screenshot_2025-12-07_at_20.57.28.png}).

\screenshot{Screenshot_2025-12-07_at_20.57.28.png}{Выбор пути для сохранения}

После успешного экспорта отображается сообщение с количеством сохранённых контактов (рис.~\ref{fig:Screenshot_2025-12-07_at_20.57.49.png}).

\screenshot{Screenshot_2025-12-07_at_20.57.49.png}{Сообщение об успешном экспорте}

\subsection{Поиск контактов}

Поиск осуществляется в реальном времени при вводе текста. На рис.~\ref{fig:Screenshot_2025-12-07_at_20.58.03.png} показан поиск по фамилии «че»~--- отображается только контакт Черкай.

\screenshot{Screenshot_2025-12-07_at_20.58.03.png}{Поиск по фамилии}

При выборе поля «Email» и вводе «el» также находится контакт Черкай, так как его email содержит подстроку «el» (рис.~\ref{fig:Screenshot_2025-12-07_at_20.58.16.png}).

\screenshot{Screenshot_2025-12-07_at_20.58.16.png}{Поиск по email}

\subsection{Сортировка таблицы}

Сортировка реализована с трёхступенчатым переключением. При первом клике по заголовку столбца «Фамилия» таблица сортируется по возрастанию (рис.~\ref{fig:Screenshot_2025-12-07_at_21.15.08.png})~--- контакты расположены в алфавитном порядке: Барашева, Иванов, Черкай. Индикатор сортировки (стрелка вверх) отображается в заголовке столбца.

\screenshot{Screenshot_2025-12-07_at_21.15.08.png}{Сортировка по фамилии по возрастанию}

При втором клике порядок меняется на убывание (рис.~\ref{fig:Screenshot_2025-12-07_at_21.15.16.png})~--- контакты расположены в обратном порядке: Черкай, Иванов, Барашева. Индикатор меняется на стрелку вниз.

\screenshot{Screenshot_2025-12-07_at_21.15.16.png}{Сортировка по фамилии по убыванию}

При третьем клике сортировка сбрасывается и восстанавливается исходный порядок добавления контактов (рис.~\ref{fig:Screenshot_2025-12-07_at_21.15.20.png}). Индикатор сортировки скрывается.

\screenshot{Screenshot_2025-12-07_at_21.15.20.png}{Сброс сортировки к исходному порядку}

\subsection{Импорт из YAML}

Для демонстрации импорта сначала очищаем справочник~--- удаляем все контакты (рис.~\ref{fig:Screenshot_2025-12-07_at_21.15.33.png}).

\screenshot{Screenshot_2025-12-07_at_21.15.33.png}{Пустой справочник перед импортом}

При нажатии кнопки «Импорт YAML» открывается диалог выбора файла (рис.~\ref{fig:Screenshot_2025-12-07_at_21.16.08.png}). Выбираем ранее экспортированный файл \texttt{phonebook.yaml}.

\screenshot{Screenshot_2025-12-07_at_21.16.08.png}{Диалог импорта из YAML}

После выбора файла появляется диалог подтверждения с количеством контактов в файле (рис.~\ref{fig:Screenshot_2025-12-07_at_21.16.25.png}).

\screenshot{Screenshot_2025-12-07_at_21.16.25.png}{Подтверждение импорта}

После подтверждения контакты загружаются в справочник и отображается сообщение об успешном импорте (рис.~\ref{fig:Screenshot_2025-12-07_at_21.16.37.png}).

\screenshot{Screenshot_2025-12-07_at_21.16.37.png}{Успешный импорт контактов}

\subsection{Формат YAML-файла}

Экспортированный файл \texttt{phonebook.yaml} имеет следующую структуру:

\begin{lstlisting}[language=C, caption=Содержимое файла phonebook.yaml]
- lastName: Черкай
  firstName: Елизаета
  middleName: ""
  address: Пушктнская д.1
  email: elizaveta.cherkay@yandex.ru
  birthDate: 2002-01-01
  phones:
    - type: mobile
      number: "79999569098"
- lastName: Иванов
  firstName: Иван
  middleName: Иванович
  address: Евграфово д.5
  email: ivanov.ii@yandex.ru
  birthDate: 2000-01-01
  phones:
    - type: mobile
      number: "79216445613"
    - type: work
      number: "79998215318"
- lastName: Барашева
  firstName: Полина
  middleName: ""
  address: Литейный д. 15
  email: p.barasheva@yandex.ru
  birthDate: 2002-01-01
  phones:
    - type: mobile
      number: "79215661718"
\end{lstlisting}

Каждый контакт представлен как элемент списка со вложенными полями. Телефоны хранятся как вложенный список с указанием типа и номера. Дата хранится в формате ISO 8601 (ГГГГ-ММ-ДД).